% Copy-paste template for TikZ diagrams
% Prevents text overlap through safe node styles and generous spacing
%
% USAGE:
%   1. Copy entire template
%   2. Replace node labels and connections
%   3. Use safe node styles: safenode, compactnode, widenode
%   4. Use safe arrow style: safearrow
%
% SAFE NODE STYLES (from tikz_diagram_style.tex):
%   safenode     - Standard node with text wrapping (2.2cm width)
%   compactnode  - Small node for short labels (1.5cm width)
%   widenode     - Wide node for long text (3cm width)
%   safearrow    - Thick blue arrow with rounded corners
%
% NODE DISTANCE:
%   Vertical:   1.8cm (generous spacing prevents overlap)
%   Horizontal: 2.5cm (prevents long labels from colliding)

\begin{figure}[htbp]
\centering

\begin{tikzpicture}[
  node distance=1.8cm and 2.5cm,  % Vertical and horizontal spacing
  every node/.style={align=center}  % Center text in nodes
]

% ============================================================
% NODES
% ============================================================

% Use safenode for standard nodes with wrapped text
\node[safenode] (data) {Training\\Data};
\node[safenode, below=of data] (features) {Feature\\Engineering};
\node[safenode, below=of features] (model) {Model\\Training};
\node[safenode, below=of model] (prediction) {Prediction};

% Use compactnode for short labels
\node[compactnode, right=of data] (validation) {Validation\\Set};

% Use widenode for long text
\node[widenode, right=of model] (ensemble) {4-Way Ensemble\\(SS + Hier + XGB + IP)};

% ============================================================
% ARROWS
% ============================================================

% Use safearrow for all connections
\draw[safearrow] (data) -- (features);
\draw[safearrow] (features) -- (model);
\draw[safearrow] (model) -- (prediction);
\draw[safearrow] (data) -- (validation);
\draw[safearrow] (model) -- (ensemble);

% Optional: Add labels to arrows
\draw[safearrow] (data) -- node[right, font=\small] {Clean} (features);

% Optional: Curved arrows (avoid overlap)
\draw[safearrow, bend right=20] (validation) to (prediction);

\end{tikzpicture}

\caption{Your caption here}
\label{fig:your-label}
\end{figure}

% ============================================================
% COMMON PATTERNS
% ============================================================

% 1. Horizontal layout (left to right):
%    \node[safenode] (a) {A};
%    \node[safenode, right=of a] (b) {B};
%    \node[safenode, right=of b] (c) {C};

% 2. Grid layout (2x2):
%    \node[safenode] (tl) {Top Left};
%    \node[safenode, right=of tl] (tr) {Top Right};
%    \node[safenode, below=of tl] (bl) {Bottom Left};
%    \node[safenode, below=of tr] (br) {Bottom Right};

% 3. Bidirectional arrows:
%    \draw[safearrow, <->] (a) -- (b);

% 4. Dashed arrows (for optional/uncertain relationships):
%    \draw[safearrow, dashed] (a) -- (b);

% 5. Multiple nodes at same level:
%    \node[safenode] (center) {Center};
%    \node[safenode, above left=of center] (tl) {Top Left};
%    \node[safenode, above right=of center] (tr) {Top Right};

% 6. Grouped nodes with box:
%    \begin{scope}[on background layer]
%    \node[draw=accent, thick, rounded corners, fit=(a)(b)(c),
%          inner sep=5mm, fill=accent!5] {};
%    \end{scope}

% 7. Self-loop (feedback):
%    \draw[safearrow] (model) edge[loop right] node[right] {Iterate} (model);

% 8. Node positioning shortcuts:
%    above=of X        % Directly above X
%    below=of X        % Directly below X
%    right=of X        % Directly right of X
%    left=of X         % Directly left of X
%    above left=of X   % Diagonal above-left
%    above right=of X  % Diagonal above-right

% 9. Custom spacing (override node distance):
%    \node[safenode, below=3cm of X] (Y) {Y};  % 3cm below instead of 1.8cm

% 10. Invisible helper nodes (for complex layouts):
%     \node[coordinate] (helper) at ($(a)!0.5!(b)$) {};  % Midpoint between a and b
%     \draw[safearrow] (a) -- (helper) -- (c);

% ============================================================
% TROUBLESHOOTING
% ============================================================

% PROBLEM: Nodes overlap
% SOLUTION: Increase node distance or use compactnode for short labels

% PROBLEM: Arrows cross nodes
% SOLUTION: Use bend right/bend left or adjust node positions

% PROBLEM: Text too long for node
% SOLUTION: Use widenode or add manual line breaks with \\

% PROBLEM: Diagram too wide for page
% SOLUTION: Reduce node distance or rotate with \begin{sidewaysfigure}

% PROBLEM: Arrows not connecting properly
% SOLUTION: Check node names match exactly (case-sensitive)
